\section{Nazwa sekcji idk}

\subsection{Techniki promptowania i generacji architektury}

Istnieje wiele opisanych w literaturze technik promptowania związanych z generowaniem architektury.

\subsubsection{Prompt Pattern Sequences}

Autorzy artykułu "A Prompt Pattern Sequence Approach to Apply Generative AI in Assisting Software Architecture Decision-making" \cite{maranhao} proponują podejście bazujące na spojrzeniu na problem z kilku różnych perspektyw. W artykule wymieniają oni pięć person które kolejno oceniają podane oczekiwania dotyczące tworzonego projektu w formie zwięzłego opisu w kilku zdaniach. Każda z person ma dokładnie zdefiniowaną rolę i sposób, w jaki ma odpowiadać na prompty użytkownika. Badanie wykazało że istnieje optymalna kolejność promptowania kolejnych utworzonych "specjalistów" i jest to:
\begin{itemize}
	\item Software Architect Persona
	\item Technical Premises
	\item Uncertain Requirements Statement
	\item Quality Attribute Question
	\item Architectural Project Context
\end{itemize}
W artykule podane są również linki do zapisanych przykładowych konwersacji z chatbotem prowadzonych za pomocą tej techniki.

\subsubsection{ADD}

W artykule "An LLM-assisted approach to designing software architectures using ADD" \cite{add} twórcy opisują proces wdrożenia LLMa w schemacie tworzenia architektury w podejściu Attribute Driven Design. Autorzy artykułu przygotowują LLMa do pracy w schemacie podając modelowi dokładny opis metody oraz tworząc dla niego personę architekta. W każdym kroku metody model tworzy artefakt z aktualnymi decyzjami architektonicznymi, co pozwala prześledzić proces myślowy tworzenia architektury, po każdym kroku stan sprawdzany jest przez człowieka do którego należą ostateczne decyzje o dalszych krokach. Autorzy zaznaczają, iż kluczowe jest podanie dokładnego schematu i opisu iteracji w wykorzystywanej metodzie, w przeciwnym przypadku model językowy gubi się.

\subsubsection{Progressive Prompting}

W artykule "Requirements are All You Need: From Requirements to Code with LLMs" \cite{wei} autor postuluje, iż LLMy mogą zostać użyte do tworzenia architektury, ale zamiast używać używać jednego promptu, można potraktować chatbota jako konsultanta z którym architekt wspólnie krok po kroku będzie ulepszał proponowaną architekturę. Proces działania zgodnie z tą metodą wygląda tak:
\begin{itemize}
	\item Wybieramy bazowego LLMa i dostosowujemy go do potrzeb generowania architektury odpowiednimi promptami
	\item Posiadając gotowego chatbota wysyłamy do niego dokumentację projektu zawierającą cele które chcemy osiągnąć i przykłady użycia
	\item Na tej podstawie chatbot przygotowuje wymagania funkcjonalne, do których użytkownik może dowolnie poprawiać
	\item W następnym kroku chatbot modeluje komponenty systemu w object-oriented design, modyfikacje i poprawki są dostępne również na tym etapie
	\item W ostatnim kroku chatbot generuje kod aplikacji i scenariusze testowe.
\end{itemize}
W artykule podane są przykłady konwersacji z przygotowanym specjalnym chatbotem opartym o ChatGPT oraz link do dotrenowanej wersji ChatGPT użytej w podanym przykładzie.

\subsubsection{ARLO}

Nazwa ARLO pochodzi od nazwy "\textbf{A}rchitecture from a set of NL \textbf{R}equirements by leveraging \textbf{L}LMs and \textbf{O}ptimization methods", metody zaproponowanej w artykule "ARLO: A Tailorable Approach for Transforming Natural Language Software Requirements into Architecture using LLMs" \cite{arlo}. Metoda ta tworzy architekturę bezpośrednio z wymagań zapisanych w języku naturalnym wykorzystując algorytmy optymalizacji programowania liniowego. ARLO na podstawie opisu wymagań tworzy listę wymagań funkcjonalnych i tworzy Design Structure Matrix, w którym decyzje architektoniczne łączone są z aspektami projektu z odpowiednio wyznaczonymi wagami, a następnie rozwiązywane jest zadanie optymalizacji całkowitoliczbowej dokonujące wyboru optymalnego w przypadku podanego projektu systemu. ARLO nie jest zadaniem bazującym jedynie na LLMie, wykorzystanie modelu jest jedynie jednym z elementów całego systemu tworzącego architekturę.

\subsubsection{DRAFT}

DRAFT \cite{draft} jest rozwiązaniem łączącym w sobie strategię RAG i fine-tuning LLMa w ramach tworzenia architektury systemu. Na początek wybierany jest model językowy, który trenowany jest za pomocą przykładowych par kontekstów i decyzji architektonicznych, z których tworzone są również prompty typu few-shot dla fine-tuningowanego LLMa. Użytkownik tworzący architekturę podaje opis swojego problemu, który na podstawie zapamiętanych par treningowych przetwarzany jest na prompt few-shot podawany do wytrenowanego LLMa który na podstawie danych treningowych oraz lepszego propmptu generuje decyzje architektoniczne.

\subsection{Metody oceny obiektywnej jakości, zgodności z wymaganiami}

\subsubsection{Przegląd literatury}

Na bardzo wczesnym etapie projektu, gdy istnieje jedynie opis architektury (np. diagramy, widoki C4, opis komponentów),
a nie ma jeszcze kodu ani alternatywnych rozwiązań do porównania, ocena jakości musi opierać się na analizie modelu architektonicznego i
scenariuszy jakościowych, a nie na pomiarach wydajności gotowego systemu. Poniżej przedstawiono trzy uzupełniające się metody,
które nadają się do takiej wczesnej analizy: scenariuszową analizę atrybutów jakości (np. ATAM/SAAM) \cite{kazman_atam_2000},
analizę modyfikowalności ALMA \cite{bengtsson_alma_2004} oraz eksperckie przeglądy architektury z użyciem list kontrolnych
\cite{cianca_eval_arch_2000,uta_design_review_checklist}.

Pierwszą grupą technik są \emph{scenariuszowe metody oceny architektury}, takie jak SAAM i ATAM,
które formułują wymagania niefunkcjonalne w postaci scenariuszy jakościowych i badają, w jaki sposób dana architektura reaguje na te scenariusze
\cite{kazman_atam_2000}. Dla pojedynczej architektury oznacza to zdefiniowanie kluczowych scenariuszy (np. skok obciążenia, awaria komponentu, typowa zmiana funkcjonalna),
przejście przez diagramy architektury krok po kroku oraz identyfikację komponentów krytycznych, punktów wrażliwości i potencjalnych ryzyk dla wydajności,
skalowalności czy niezawodności, jeszcze zanim powstanie kod \cite{cianca_eval_arch_2000}.

Drugą metodą, wyspecjalizowaną w analizie modyfikowalności na poziomie architektury, jest
\emph{Architecture-Level Modifiability Analysis (ALMA)} \cite{bengtsson_alma_2004}.
ALMA składa się z pięciu kroków: wyboru celu analizy, doprecyzowania opisów architektury, elicytacji scenariuszy zmian,
ich systematycznej oceny oraz interpretacji wyników, co pozwala oszacować koszt oraz ryzyko typowych zmian już na etapie projektu architektury.
Ponieważ cała analiza opiera się na scenariuszach zmian i strukturze architektury, metoda dobrze nadaje się do wczesnej fazy,
gdy jeszcze nie istnieje implementacja, a celem jest wykrycie potencjalnie kosztownych obszarów (wysokie sprzężenie, wąskie gardła odpowiedzialności itp.)
\cite{bengtsson_alma_2004}.

Trzecią klasą podejść są \emph{eksperckie przeglądy architektury z użyciem list kontrolnych}
(architecture/design review checklists),
w których doświadczeni architekci przechodzą przez ustrukturyzowany zestaw pytań dotyczących m.in. prostoty i spójności podziału na komponenty,
powiązań z wymaganiami, przejrzystości interfejsów, zarządzania błędami czy aspektów bezpieczeństwa \cite{uta_design_review_checklist}.
Tego typu checklisty pozwalają zidentyfikować błędy koncepcyjne (zbyt złożone komponenty, nadmierne sprzężenie, brak jawnie zdefiniowanych interfejsów,
nieprzemyślana redundancja) zanim jakiekolwiek decyzje zostaną „utrwalone” w kodzie, a jednocześnie są stosunkowo lekkie procesowo i
dobrze łączą się z podejściami scenariuszowymi \cite{cianca_eval_arch_2000,uta_design_review_checklist}.

Podsumowując, we wczesnej fazie projektu -- kiedy dostępna jest tylko dokumentacja architektury -- sensowne jest połączenie dwóch perspektyw: 
(1) scenariuszowej analizy atrybutów jakości (np. w duchu ATAM), oraz (2) ukierunkowanej analizy modyfikowalności ALMA .
Razem pozwalają one ocenić, na ile dana architektura rokuje pod kątem kluczowych atrybutów jakości, zanim powstanie implementacja.

\subsubsection{Miarodajne metryki}

Oprócz metod scenariuszowych i przeglądów eksperckich, jakość architektury na wczesnym etapie można próbować oceniać za pomocą prostych metryk.
W literaturze dotyczącej stabilności i modyfikowalności architektury podkreśla się znaczenie miar sprzężenia,
rozmiaru oraz rozkładu odpowiedzialności jako głównych predyktorów łatwości zmian i ewolucji systemu \cite{bengtsson_alma_2004,add}.
Poniżej przedstawiono trzy przykładowe metryki.

\subsubsection{Średni wpływ zmiany (Average Change Impact)}

Pierwszą metryką jest \emph{średni wpływ zmiany}, która przybliża poziom sprzężenia między komponentami oraz modyfikowalność architektury.
Dla zadanego zestawu scenariuszy zmian (np. dodanie nowej funkcjonalności, zmiana sposobu autoryzacji, podmiana systemu płatności)
dla każdego scenariusza identyfikuje się komponent wyjściowy, w którym inicjowana jest zmiana, a następnie liczbę komponentów,
które muszą zostać zmodyfikowane, aby scenariusz zrealizować.
Metryka jest obliczana jako średnia liczba komponentów dotkniętych zmianą w przeliczeniu na scenariusz:
\[
\text{ACI} = \frac{1}{C}\frac{1}{N}\sum_{i=1}^{N} c_i,
\]
gdzie \(N\) to liczba scenariuszy, \(C\) to liczba komponentów, a \(c_i\) to liczba komponentów wymagających modyfikacji dla scenariusza \(i\).
Im niższa wartość ACI, tym architektura jest potencjalnie mniej sprzężona.

\subsubsection{Wskaźnik stabilności architektury}

Drugą metryką jest \emph{wskaźnik stabilności architektury}. 
W kontekście modeli architektonicznych można go zdefiniować jako udział komponentów,
które pozostają niezmienione w typowych scenariuszach ewolucji systemu (np. dodanie nowych przypadków użycia, integracja z zewnętrzną usługą, zmiana technologii bazy danych).
Dla przyjętego zestawu scenariuszy określa się, które komponenty wymagają modyfikacji, a które pozostają stabilne,
następnie wyznacza się stosunek liczby komponentów stabilnych do całkowitej liczby komponentów w architekturze:
\[
\text{Stab} = \sum_{i=1}^{N} \frac{|C_{\text{stabilne}}|}{|C_{\text{wszystkie}}|}.
\]
Wysoka wartość wskaźnika sugeruje, że architektura dobrze lokalizuje zmiany w ograniczonej liczbie komponentów,
co redukuje ryzyko nieoczekiwanych efektów ubocznych podczas rozwoju systemu \cite{bengtsson_alma_2004}.

\subsubsection{Średni stopień zależności komponentów}

Trzecią metryką jest \emph{średni stopień zależności komponentów}, która mierzy,
jak silnie każdy komponent jest powiązany z pozostałymi, i tym samym w jakim stopniu architektura sprzyja niskosprzężonemu, modułowemu projektowi.
Dla każdego komponentu liczy się liczbę zależności wychodzących (np. wywołań, komunikatów, relacji wymagań) do innych komponentów, a następnie oblicza średnią po wszystkich komponentach:
\[
\text{DepAvg} = \frac{1}{|C|}\sum_{k \in C} d_k,
\]
gdzie \(C\) to zbiór komponentów, a \(d_k\) to liczba zależności wychodzących komponentu \(k\). Można dodatkowo rozróżniać zależności krytyczne (np. w ścieżkach transakcji biznesowych) i wagować je wyżej w obliczeniach. Niższa wartość tej metryki wskazuje na mniejszy poziom sprzężenia i potencjalnie lepszą skalowalność oraz testowalność architektury, co jest spójne z obserwacjami dotyczącymi roli sprzężenia i granic komponentów w jakości projektów architektonicznych \cite{add}.

\section{Dane testowe}

W celu przeprowadzenia eksperymentów przeanalizowano publicznie dostępne aplikacje otwartoźródłowe oparte na architekturze mikroserwisowej. Analiza skupiła się na weryfikacji dostępnych artefaktów: dokumentacji architektonicznej (typ i format diagramów, modelowane aspekty systemu, liczba zdefiniowanych serwisów) oraz wymaganiach systemowych (obecność historyjek użytkownika (ang. \emph{user stories}), stopień ustrukturyzowania, podział na wymagania funkcjonalne i niefunkcjonalne).

W pierwszej kolejności wyselekcjonowano zbiory posiadające zadowalająco opisane zarówno wymagania, jak i docelową dokumentację architektoniczną. Kluczowym źródłem danych w tej kategorii jest praca \textit{A curated Dataset of Microservices-Based Systems} \cite{imranur_curated_2019}, w której autorzy zagregowali i opisali 20 projektów otwartoźródłowych. W obrębie tego zbioru na szczególną uwagę zasługują następujące aplikacje:

\paragraph{eShop (eShopOnContainers)} \cite{noauthor_dotneteshop_2026} -- rozbudowana aplikacja referencyjna dla platformy .NET, składająca się z wielu kooperujących mikroserwisów (m.in. koszyk, katalog, tożsamość użytkownika, zamówienia). Dostarcza wysoce ustrukturyzowane wymagania w formie wypunktowanych list, z wyraźnym podziałem na wymagania funkcjonalne (np. zarządzanie koszykiem) i niefunkcjonalne (wymagana wysoka dostępność, wsparcie dla wielu klientów front-endowych). Dokumentacja zawiera szczegółowe wizualne diagramy architektoniczne, obrazują one strukturę systemu, mapowanie aplikacji klienckich przez API Gateway do serwisów docelowych oraz infrastrukturę komunikacji asynchronicznej opartej na szynie zdarzeń (ang. \textit{Event Bus}).
\paragraph{Pitstop} \cite{wijk_edwinvwpitstop_2026} -- aplikacja do zarządzania warsztatem samochodowym obrazująca architekturę opartą na wzorcach Event Sourcing oraz CQRS i zaprojektowana zgodnie z metodyką BDD. Projekt nie posiada spisanych historyjek użytkownika, ale dostarcza bogaty i sformalizowany model domeny oraz wyniki warsztatów \textit{Event Storming}. Architektura dokumentowana jest diagramami ilustrującymi mapę kontekstów (ang. \textit{Context Map}) oraz schematem przepływu poszczególnych komend i zdarzeń pomiędzy serwisami za pomocą brokera wiadomości.
\paragraph{Spring Petclinic Microservices} \cite{noauthor_spring-petclinicspring-petclinic-microservices_2026} -- rozproszona wersja klasycznej aplikacji złożona z kilku głównych domenowych mikroserwisów i infrastruktury bazującej na Spring Cloud. Projekt zawiera zwięzły opis domenowy (oparty na zarządzaniu weterynarzami oraz wizytami), z kolei wymagania niefunkcjonalne nie zostały formalnie wyodrębnione. Zawarto w nim także nieformalne scenariusze użycia pod postacią zapytań do chatbota opartego na systemach GenAI. Architektura topologiczna jest dobrze udokumentowana za pomocą czytelnego diagramu załączonego jako statyczny obraz w głównym pliku dokumentacji.

Dodatkowo do zbiorów bogatych w artefakty zaliczono systemy badane w pracy \cite{slater_quantitative_2025}:

\paragraph{EventTicketSystem oraz HotelPricingSystem} -- aplikacje te dostarczają bardzo jakościowy zestaw zarówno wymagań funkcjonalnych, jak i niefunkcjonalnych. Odtworzone dla nich struktury wyjściowe stanowią obszerną bazę badawczą i architektoniczną w postaci dynamicznych diagramów (stworzonych w języku skryptowym Mermaid) na różnych poziomach abstrakcji oraz szczegółowo zdefiniowanych, ustrukturyzowanych zapisów decyzji architektonicznych (ang. \textit{Architecture Decision Records} -- ADR).

Dla kontrastu, przeanalizowano również zbiory dostarczające rozbudowaną architekturę bazową, w których brakuje jednak ustrukturyzowanych wymagań systemowych, które mogłyby stanowić wejście dla modeli generatywnych:

\paragraph{Train Ticket} \cite{noauthor_fudanselabtrain-ticket_2026} -- rozbudowany projekt składający się z 41 mikroserwisów. W jego repozytorium brakuje formalnych i nieformalnych założeń biznesowych.
\paragraph{Online Boutique} \cite{noauthor_googlecloudplatformmicroservices-demo_2026} -- demonstracyjna aplikacja sklepu (złożona z 10 usług), skupiona wokół demonstracji interakcji infrastrukturalnych w środowisku Kubernetes i technologii gRPC.

\section{Podejścia do ewaluacji wyników}

Weryfikacja docelowej architektury wygenerowanej na podstawie wymagań poprzez wielkie modele językowe wymaga przyjęcia powtarzalnych i miarodajnych metod oceny. Na podstawie analizowanej literatury wyróżnia się następujące techniki ewaluacyjne:

\paragraph{Porównanie z referencyjnym grafem komunikacji (ang. \textit{Reference Service Graph})} -- metoda opisana w pracy \textit{Toward Generating Microservice Architectures from Textual Requirements with Large Language Models} \cite{pereiraGeneratingMicroserviceArchitectures2025}. Podejście to polega na bezpośrednim zestawieniu grafu wygenerowanych zależności pomiędzy usługami ze wzorcową architekturą weryfikacyjną (ang. \textit{ground truth}). Na tej podstawie dokonywane jest ilościowe wyliczenie metryk takich jak trafność (ang. \textit{precision}), czułość (ang. \textit{recall}) oraz miara F1 na poziomie odzyskanych mikroserwisów i połączeń między nimi. Służy to badaniu stopnia poprawnej dekompozycji mająd do dyspozycji architekturę ,,referencyjną''.
\paragraph{Ekstrakcja faktów z wymagań architektonicznych} -- koncepcja zaproponowana w pracy \textit{Do Large Language Models Contain Software Architectural Knowledge?} \cite{soliman_large_2025}. Polega na identyfikacji określonych faktów projektowych i zależności ze wstępnych wymagań (np. "moduł X wykonuje operację Y", "usługa X asynchronicznie przesyła dane do usługi Z"), a następnie zweryfikowaniu, czy system zrekonstruowany przez model językowy fizycznie włącza i implementuje wspomniane założenia. Gwarantuje to ocenę zgodności otrzymanego rozwiązania z narzuconymi z góry atrybutami poprawności i bezpieczeństwa.
\paragraph{Zastosowanie obliczalnych metryk architektonicznych} -- przegląd systematyczny \textit{Software Architecture Metrics: a literature review} \cite{coulin_software_2019} dowodzi skuteczności użycia wyliczalnych, heurystycznych miar oceniających graf systemu. W odniesieniu do architektury mikroserwisowej kluczowe okazuje się badanie poziomu powiązań pomiędzy serwisami (ang. \textit{coupling}), miary ich wenętrznej spójności (ang. \textit{cohesion}) oraz ogólnego narzutu złożoności wynikowej aplikacji niezależnie od architektury referencyjnej. Niskie sprzężenie oraz wysoka spójność to ugruntowane inżynieryjnie metryki poprawności projektu.
\paragraph{Wykorzystanie miar z dziedziny przetwarzania języka naturalnego (NLP)} -- podejście zaproponowane w pracy \textit{ArchBench: Benchmarking Generative-AI for Software Architecture Tasks} \cite{adnan_archbench_2026} w kontekście zadań takich jak generowanie zapisów decyzji architektonicznych (ADR). Zastosowanie klasycznych metryk oceny tekstu (ROUGE, BLEU, METEOR, BERTScore) pozwala na pomiar różnic, precyzji oraz semantycznego podobieństwa pomiędzy wygenerowanymi a referencyjnymi opisami architektury w postaci czystego tekstu (ang. \textit{plaintext}), choć wyjaśnialność takich wyników jest mniej oczywista.

